\documentclass[a4paper]{article}     
\usepackage{amsfonts}
\usepackage{amsmath}
\usepackage{amssymb}
\usepackage{graphicx}
\usepackage{color}

\newcommand{\Q}{\mathbb{Q}}
\newcommand{\R}{\mathbb{R}}
\newcommand{\llim}{\lim\limits}
\newcommand{\dint}{\displaystyle\int}
\newcommand{\Bgtan}{\operatorname{Bgtan}} 

\begin{document}

\noindent \large Auteur: Hina Rana \\
\noindent \large Studierichting: Informatica\\
\noindent \large Oefeningenreeks 2B nr. 1.2\\

\medskip

\normalsize

$\sin (Bgtanx)$
neem $y= \Bgtan(x)$ met $ y \in  \left]-\dfrac{\pi}{2},\dfrac{\pi}{2}\right[\\

$\tan(y) =x$\\

$\tan y \cdot \cos (y) = x\cdot\cos(y)$\\

$\sin y = x\cos(y)$\\

$\sin^2 y + \cos^2(y) =1$\\

$(x\cos y)^2 + \cos^2(y)=1$\\

$x^2\cos^2y + \cos^2(y) =1$\\

$\cos^2 y(x^2+1) =1$\\

$\cos^2 y = \dfrac{1}{x^2+1}$\\


$\cos y = \dfrac{1}{\sqrt{x^2+1}}$\\


$x\cos y = \dfrac{x}{\sqrt{x^2+1}}\\


$\sin y = \dfrac{x}{\sqrt{x^2+1}}$\\

\begin{thebibliography}{999}
\end{thebibliography}
\end{document} 
